\documentclass[UTF8,landscape,12pt]{ctexart}
\usepackage[a3paper,landscape,margin=1.2cm]{geometry}
\usepackage{graphicx}
\usepackage{xcolor}
\usepackage{amsmath}
\usepackage{enumitem}
\usepackage{tcolorbox}
\usepackage{multicol}
\usepackage{caption}
\pagestyle{empty}
\graphicspath{{assets/}}

\definecolor{TsinghuaPurple}{HTML}{6A1B78}
\definecolor{DeepPurple}{HTML}{3C1053}
\definecolor{LightPurple}{HTML}{F5EDF7}
\definecolor{Ink}{HTML}{222222}
\definecolor{Muted}{HTML}{666666}
\captionsetup{font=footnotesize,labelfont=bf}
\setlength{\parindent}{0pt}
\setlist[itemize]{leftmargin=1.2em,itemsep=0.12em,topsep=0.1em}

\newtcolorbox{panel}[1]{
  colback=white,
  colframe=LightPurple,
  boxrule=0.8pt,
  arc=2pt,
  left=5pt,right=5pt,top=4pt,bottom=4pt,
  title=\bfseries #1,
  coltitle=TsinghuaPurple,
  colbacktitle=LightPurple,
  fonttitle=\small
}

\begin{document}
\color{Ink}

\begin{minipage}[t]{0.72\textwidth}
  {\fontsize{24}{28}\selectfont\bfseries\color{TsinghuaPurple} 表面重力波波长与表观波速的水深效应探究}\\[0.2em]
  {\large 探究性实验成果海报}
\end{minipage}
\hfill
\begin{minipage}[t]{0.25\textwidth}
  \raggedleft
  \includegraphics[width=0.55\linewidth]{tsinghua_logo.png}\\[-0.2em]
  {\small 丘洪源、白子恒、秦健淞、纪文龙、杨天皓\\行健书院\quad 2026 年 5 月}
\end{minipage}

\vspace{0.35cm}
\colorbox{LightPurple}{%
  \begin{minipage}{0.98\textwidth}
    \vspace{0.25em}
    \textbf{核心问题：} 水深是否会明显改变水波传播？在实际水槽中，我们能严格测到理论相速度，还是只能得到受反射影响的表观波速？
    \vspace{0.25em}
  \end{minipage}}

\vspace{0.35cm}
\begin{minipage}[t]{0.31\textwidth}
  \begin{panel}{实验设计}
    \begin{itemize}
      \item 透明水槽、2 Hz 活塞电机、3D 打印推板、LED 背光和侧向摄像。
      \item 水深设置为 $2.2$--$20.0\,\mathrm{cm}$ 共 14 组。
      \item 从视频中重建 $\eta(x,t)$，提取推荐平均波长 $\lambda_{\mathrm{rec}}$。
      \item 追踪波峰轨迹，统计中位表观速度 $|c_{\mathrm{app}}|$ 及离散范围。
    \end{itemize}
    \centering
    \includegraphics[width=0.95\linewidth]{fig_setup.png}\\[-0.2em]
    {\footnotesize 图 1\quad 侧向成像与液面识别}
  \end{panel}

  \vspace{0.25cm}
  \begin{panel}{关键定义}
    \[
      \omega^2=\left(gk+\frac{\sigma}{\rho}k^3\right)\tanh(kh),\qquad
      c_p=\frac{\omega}{k}
    \]
    \small
    本实验报告的速度是视频追踪得到的 $|c_{\mathrm{app}}|$，不是严格理论相速度 $c_p$。这是解释结果时最重要的边界。
  \end{panel}
\end{minipage}
\hfill
\begin{minipage}[t]{0.34\textwidth}
  \begin{panel}{视频证据}
    \centering
    \includegraphics[width=0.98\linewidth]{fig_xt_field.png}\\[-0.2em]
    {\footnotesize 图 2\quad $h=12\,\mathrm{cm}$ 的 $x$--$t$ 位移场与传播方向分离}
    \vspace{0.25em}
    \small
    连续斜纹说明水面波列可以被稳定追踪；方向分离显示入射波和反射波同时存在。
  \end{panel}

  \vspace{0.25cm}
  \begin{panel}{理论判据}
    \centering
    \includegraphics[width=0.98\linewidth]{fig_theory_compare.png}\\[-0.2em]
  \end{panel}
\end{minipage}
\hfill
\begin{minipage}[t]{0.31\textwidth}
  \begin{panel}{主要结果}
    \centering
    \includegraphics[width=0.98\linewidth]{fig_lambda_capp.png}\\[-0.2em]
    {\footnotesize 图 4\quad 波长与表观波速随水深变化}
    \vspace{0.25em}
    \small
    除最低水深组外，$\lambda_{\mathrm{rec}}$ 主要集中在 $8\,\mathrm{cm}$ 附近；表观波速离散更大，没有简单单调律。
  \end{panel}

  \vspace{0.25cm}
  \begin{panel}{结论与改进}
    \small
    \begin{itemize}
      \item 多数组按 $kh$ 判据已接近深水近似区，水深修正本身较弱。
      \item 表观波速受反射叠加和波峰追踪稳定性影响，不能直接等同于理论相速度。
      \item 后续应加入消波装置，并同步提取主频 $f$，用 $c=\lambda f$ 与理论相速度比较。
    \end{itemize}
  \end{panel}
\end{minipage}

\end{document}
